\section{Nonlinear optics equation on a semi-strip} \label{NW}

In this section we consider the nonlinear optics equation \eqref{NW1}, where $D$ has the form \eqref{NW2}.
We consider equation \eqref{NW1} on the semi-strip $\cD$, which is given by~\eqref{1.4}. First, using Weyl theoretic
results from Section~\ref{subs2.2}, we express Weyl function $\vp(t,z)$ of the auxiliary system~\eqref{NW5},
where $\zeta(x)=\zeta(x,t)=[D,\vr(x,t)]$, in terms of $\vp(0,z)$ and
the boundary condition $\vr(0,t)=\wh \rho(t)$. In other words, we express in these terms the evolution of the Weyl function.
For that purpose, following formulas~\eqref{2.4} and~\eqref{2.5} in Proposition~\ref{TmM},
we introduce matrix function $R(t,z)$ by the relations
\begin{gather} \label{NW20}
R_t(t,z)=(\I z \wh D-[\wh D,\wh \rho(t)])R(t,z), \qquad R(0,z)=I_m.
\end{gather}
Recall that $\wh D$ is a diagonal matrix and that $\wh D>0$.
\begin{Theorem}
\label{TmNWevol}
Let $\vr(x,t)$ satisfy the nonlinear optics equation \eqref{NW1} and the boundary condition $\vr(0,t)=\wh \rho(t)$.
Assume that $\vr$ is uniformly bounded and continuously differentiable on~$\cD$.
Then, the matrix functions
\begin{gather} \label{NW21}
\breve \psi_k(t,z):=
 \begin{bmatrix} I_k & 0\end{bmatrix}
R(t,z) \varphi (0,z)
 \left[
  \begin{matrix}
    0  \\
    I_{m-k}
  \end{matrix}
 \right]
 \biggl(
 \begin{bmatrix} 0 & I_{m-k}\end{bmatrix}
  R(t,z) \varphi (0,z)
  \left[
  \begin{matrix}
    0  \\
    I_{m-k}
  \end{matrix}
  \right]
 \biggr)^{-1}
\end{gather}
are well-defined for $1\leq k <m$, and the evolution of the Weyl function is given by the formula
\begin{gather} \label{NW22}
  \left\{
   \varphi_{i, k+1}(t,z)
  \right\}_{i=1}^k
 = \breve \psi_k (t,z)\begin{bmatrix} 1\\ 0 \\ \ldots \\0
 \end{bmatrix}, \qquad \Im(z)<-M
\end{gather}
and by the normalization conditions~\eqref{NWd7}.
\end{Theorem}

\begin{proof} We set
\begin{gather} \label{NW23}
M_0=\sup\| \zeta(x,t)\|, \qquad (x,t)\in \cD, \qquad M>\max_{1\leq k<m}\big(2M_0/(d_k-d_{k+1})\big).
\end{gather}
Recall that $G$ and $F$ for the case of the nonlinear optics equation are given by~\eqref{NW4}. Since $\vr$ is continuously dif\/ferentiable, the conditions of Proposition~\ref{TmM}
are fulf\/illed. Taking into account that the fundamental solution of~\eqref{NW5} is denoted by~$w$
(instead of~$u$ in Proposition~\ref{TmM}), we rewrite~\eqref{2.6} in the form
$w(x,t,z)=R(x,t,z)w(x,0,z)R(t,z)^{-1}$ or, equivalently,
\begin{gather} \label{NW24}
w(x,t,z)^{-1}=R(t,z)w(x,0,z)^{-1}R(x,t,z)^{-1}.
\end{gather}
Using \eqref{NW24}, we express $w(x,t,z)^{-1}$ via $w(x,\wt t,z)^{-1}$, $0 \leq t,\wt t <a$:
\begin{gather}
w(x,t,z)^{-1} =R(t,z)R(\wt t,z)^{-1}R(\wt t,z)w(x,0,z)^{-1}R(x,\wt t,z)^{-1}\big(R(x,t,z)R(x,\wt t,z)^{-1}\big)^{-1}\nonumber\\
 \label{NW25}
 \hphantom{w(x,t,z)^{-1}}{}
 =R(t,z)R(\wt t,z)^{-1}w(x,\wt t,z)^{-1}\big(R(x,t,z)R(x,\wt t,z)^{-1}\big)^{-1}.
\end{gather}
Recall that $R_t=FR$, where $F$ is given in \eqref{NW4} and $\wh D>0$. Hence, putting $\cR(x,t, \wt t, z):=\big(R(x,t,z)R(x,\wt t,z)^{-1}\big)^{-1}$ we derive
\begin{gather*}
\frac{\prt}{\prt t} \cR(x,t, \wt t, z)=-\cR(x,t, \wt t, z)F(x,t,z), \\
\frac{\prt}{\prt t} \big(\cR(x,t, \wt t, z)\cR(x,t, \wt t, z)^*)<0, \qquad \Im(z)<0 , \qquad
 \cR(x,\wt t, \wt t, z)=I_m.
\end{gather*}
From the relations above it is immediate that
\begin{gather} \label{NW26}
\|\cR(x, t, \wt t, z)\|<1 \qquad {\mathrm{for}}\quad t>\wt t, \qquad \Im(z)<0,
\end{gather}
which allows us to estimate the dif\/ference
\begin{gather} \label{NW27}
I_m-\cR(x, t, \wt t, z)=\int_{\wt t}^t \cR(x, s, \wt t, z)F(x,s,z)ds.
\end{gather}
According to \eqref{NW23}, \eqref{NW26} and \eqref{NW27}, for each $\delta >0$ and $c>M$ there is $\ve=\ve(M_0)>0$ such that
\begin{gather} \label{NW28}
\| I_m-\cR(x, t, \wt t, z)\| \leq \delta
\\ \nonumber \text{for all} \quad x\in [0, \infty), \qquad \! 0\leq t-\wt t \leq \ve,\qquad\!
 z \in \{z \colon |z|<c\}\cap \{z \colon \Im(z)<-M\}, \qquad\! c>M .
\end{gather}
Modifying~\eqref{NW13} (so that the functions $\psi_k$ and $w$ depend there also on an additional variable~$t$),
in view of~\eqref{NW25},
we derive
\begin{gather}\nonumber
  \psi_k (x,t,z)
= \begin{bmatrix}I_k &0 \end{bmatrix}
R(t,z)R(\wt t,z)^{-1}w(x,\wt t,z)^{-1}\cR(x,t, \wt t,z)
 \cQ_k (z)
\\
\hphantom{\psi_k (x,t,z)=}{} \times \left(
  \begin{bmatrix}0 &I_{m-k} \end{bmatrix}
 R(t,z)R(\wt t,z)^{-1}w(x,\wt t,z)^{-1}\cR(x,t, \wt t,z)
  \cQ_k(z)
 \right)^{-1}.
 \label{NW30}
\end{gather}
Moreover, putting
\begin{gather} \label{NW31}
\wt \cQ_k(z):=\cR(x,t, \wt t,z)\cQ_k(z), \qquad \cQ_k(z):=\begin{bmatrix} 0 & I_{m-k}\end{bmatrix}
\end{gather}
we see that for suf\/f\/iciently small $\delta$ the matrix function $\wt \cQ_k(z)$ satisf\/ies \eqref{NW12} in the domain (for~$z$)
given in~\eqref{NW28}. Substituting $\wt \cQ_k$ (instead of $ \cQ_k$) into~\eqref{NW13}, we obtain
\begin{gather} \nonumber
\begin{bmatrix}\psi_k(x, \wt t,z) \\ I_{m-k} \end{bmatrix}=w(x,\wt t,z)^{-1}\wt \cQ_k(z)\left(
  \begin{bmatrix}0 &I_{m-k} \end{bmatrix}
 w(x,\wt t,z)^{-1}
  \wt \cQ_k(z)
 \right)^{-1}, \qquad {\mathrm{i.e.,}}
\\ \label{NW32}
w(x,\wt t,z)^{-1}\wt \cQ_k(z)=\begin{bmatrix}\psi_k(x, \wt t,z) \\ I_{m-k} \end{bmatrix}\left(
  \begin{bmatrix}0 &I_{m-k} \end{bmatrix}
 w(x,\wt t,z)^{-1}
  \wt \cQ_k(z)
 \right).
\end{gather}
Using the f\/irst equality in \eqref{NW31}, we can substitute \eqref{NW32} into \eqref{NW30}.
Thus, we derive
\begin{gather*}
\begin{bmatrix} \psi_k(x, t,z) \\ I_{m-k} \end{bmatrix}
=
R(t,z)R(\wt t,z)^{-1}\begin{bmatrix}\psi_k(x, \wt t,z) \\ I_{m-k} \end{bmatrix} \left(
  \begin{bmatrix}0 &I_{m-k} \end{bmatrix}
 R(t,z)R(\wt t,z)^{-1}\begin{bmatrix}\psi_k(x, \wt t,z) \\ I_{m-k} \end{bmatrix} \right)^{-1}
\end{gather*}
and in view of \eqref{NW15}, passing to the limit $x \to \infty$, we have the following formula{\samepage
\begin{gather}\nonumber
\begin{bmatrix}\breve \psi_k( t,z) \\ I_{m-k} \end{bmatrix}
=
R(t,z)R(\wt t,z)^{-1}\begin{bmatrix}\breve \psi_k( \wt t,z) \\ I_{m-k} \end{bmatrix}
\\ \label{NW33}
\hphantom{\begin{bmatrix}\breve \psi_k( t,z) \\ I_{m-k} \end{bmatrix}
= }{}\times
 \left(
  \begin{bmatrix}0 &I_{m-k} \end{bmatrix}
 R(t,z)R(\wt t,z)^{-1}\begin{bmatrix}\breve \psi_k(\wt t,z) \\ I_{m-k} \end{bmatrix} \right)^{-1}.
\end{gather}}

\noindent
Here we used the fact that, according to \eqref{NW14} and \eqref{NW28},
\begin{gather*}
\det \left(
  \begin{bmatrix}0 &I_{m-k} \end{bmatrix}
 R(t,z)R(\wt t,z)^{-1}\begin{bmatrix}\breve \psi_k(\wt t,z) \\ I_{m-k} \end{bmatrix} \right)\not=0
\end{gather*}
for suf\/f\/iciently small $\delta$. Setting $\wt t=k\ve \,\, (k=0,1, \ldots)$, recalling that $R(0,z)=I_m$ and applying each time~\eqref{NW33}, we easily prove (by induction) the equality
\begin{gather} \label{NW34}
\begin{bmatrix}\breve \psi_k( t,z) \\ I_{m-k} \end{bmatrix}
=
R(t,z)\begin{bmatrix}\breve \psi_k(0,z) \\ I_{m-k} \end{bmatrix}
 \left(
  \begin{bmatrix}0 &I_{m-k} \end{bmatrix}
 R(t,z)\begin{bmatrix}\breve \psi_k(0,z) \\ I_{m-k} \end{bmatrix} \right)^{-1}
\end{gather}
for $t$ on all intervals $[0, (k+1) \ve]\cap [0,a)$, that is, for $t$ on $[0,a)$. Although
\eqref{NW34} is proved for $ z \in \{z \colon |z|<c\}\cap \{z \colon \Im(z)<-M\}$, the analyticity
of both sides of \eqref{NW34} implies that the equality holds in the semi-plane $\Im(z)<-M$.
Finally, we note that~\eqref{NW17} at $t=0$ yields
\begin{gather} \label{NW35}
\begin{bmatrix}\breve \psi_k( 0,z) \\ I_{m-k} \end{bmatrix}
=
\vp(0,z)\begin{bmatrix}0 \\ I_{m-k} \end{bmatrix}
 \left(
  \begin{bmatrix}0 &I_{m-k} \end{bmatrix}
 \vp(0,z)\begin{bmatrix} 0 \\ I_{m-k} \end{bmatrix} \right)^{-1}.
\end{gather}
Substituting \eqref{NW35} into~\eqref{NW34}, we obtain~\eqref{NW21}. The procedure to recover
$\vp(t,z)$ from $\{\breve \psi_k(t,z)\}$ (for f\/ixed values of $t$) is described in Section~\ref{subs2.2}
(note that~\eqref{NW22} coincides with~\eqref{NW16}).
\end{proof}
